\chapter{Why synthetic data, and what it costs}
\label{ch:synthetic}

\begin{saybox}[title=What to say at the start of this section]
This is the section most likely to be interrupted with a direct question.
Do not begin with a justification. Begin with the honest statement of what
the repository can and cannot show, then give the scientific argument for
why synthetic expansion was an intelligible choice, then state plainly
what it costs. The panel will trust the argument more if the limitation
is volunteered rather than extracted.
\end{saybox}

\section{The evidence limitation, stated first}
\label{sec:limit-evidence}

Before any argument about motivation, the following must be stated,
because it constrains everything that can honestly be claimed afterwards.

\begin{limitbox}[title=Limitation 1 --- the original experimental observations are not in this repository]
Every dataset file used for training at any recoverable point in this
project's history was loaded and its \texttt{data\_origin} column
value-counted. The result appears in \Cref{tab:tab_dataset_versions}: across
five successive dataset versions spanning 30\,500 to 48\,240 rows, every
single \texttt{data\_origin} value is a synthetic-generation tag. No file
contains a \texttt{real}, \texttt{experimental}, or
\texttt{real\_paper\_derived} row.

Consequently the original experimental or literature-extracted
observations that must have calibrated the generation process cannot be
counted, characterised, or audited from this repository. Any figure for
their number --- including the figure of roughly 400 that has been quoted
informally --- is unverifiable here. This report does not assert one.
\end{limitbox}

\begin{limitbox}[title=Limitation 2 --- the generation code is not in this repository]
Every Python file in the repository was searched for code that
\emph{assigns} \texttt{source\_rule\_id}, as opposed to reading it for
grouping or splitting. There is none. Every script that touches the column
--- the training driver, the serving layer, the diagnostic experiments ---
consumes it as a pre-existing column of an already-generated file. The
parameter distributions, sampling ranges, target equations, noise model and
censoring mechanism of the generator are therefore not inspectable here.

Everything this report says about generation-rule behaviour is inferred
\emph{statistically from the generated data} and is labelled as such. No
generation equation is reconstructed or guessed.
\end{limitbox}

\tbl{tab_dataset_versions}{Every training dataset version in the
repository, with the \texttt{data\_origin} values actually present in it.
Produced by loading each file and value-counting the column.}

\noindent A direct consequence for the meeting: the honest answer to
``how many real observations did you start with'' is that the number is
not recoverable from this repository and should be requested from whoever
maintains the generation pipeline. Stating that is stronger than quoting a
number that cannot be substantiated.

\section{Why a small experimental corpus is hard to model, stated precisely}

The claim ``there was not enough data'' is too loose to defend in front of
a panel, and it is also not quite the right claim. There is no universal
minimum row count for a Random Forest. What matters is the relationship
between the diversity of the relationships to be learned and the number of
\emph{independent} observations supporting them. Four specific structural
difficulties apply here, and each can be quantified from the data that
does exist.

\subsection{Fragmentation across specialists}

The architecture divides whatever data exists by six. \Cref{fig:fig04_dataset_composition}
shows the division for the synthetic corpus, and the same division would
apply to any real corpus. If a real corpus of a few hundred observations
were distributed across six specialists in the proportions the synthetic
corpus uses, the smallest specialist --- hard/safety-endpoint, holding
$6.6\%$ of the corpus --- would receive on the order of tens of rows.
Tens of rows cannot support the estimation of a nonlinear function over
twenty-plus heterogeneous inputs, and this conclusion does not depend on
knowing the exact original count.

\subsection{Dimensionality and heterogeneity of the input space}

The feature space is not merely high-dimensional; it is
\emph{combinatorial}. A single prediction is conditioned on the cheese
product, its physical form, its composition, the treatment and its
concentration, the application method, the storage temperature, the
packaging and its headspace gas composition, and the specific indicator
being tracked with its threshold. In the V7 corpus the hard/general
specialist alone spans six base cheese products, four treatment
ingredients including the untreated control, and eight distinct indicator
types --- and that is after synthetic expansion. The product of the
observed level counts is far larger than any realistic experimental
programme could populate.

\subsection{Sparse and uneven coverage}

Coverage is not uniform across that space, and the unevenness is severe
even in the expanded corpus. \Cref{fig:fig06_rule_size_distribution,fig:fig07_rule_concentration}
show it directly: in the soft/general specialist, one generation rule
contributes $5\,654$ of $8\,800$ rows while two others contribute $17$ and
$8$ rows respectively --- a ratio of roughly $700$:$1$ between the largest
and smallest. Regions of the space represented by single-digit row counts
cannot support reliable estimation, and \Cref{ch:folds} shows exactly what
happens when the model is asked to predict them.

\subsection{Statistical estimability}

The relevant formal statement is the bias--variance decomposition. For a
model $\hat{f}$ trained on a random sample $\mathcal{D}$, the expected
squared error at a point $\mathbf{x}_0$ decomposes as
\begin{equation}
\mathbb{E}_{\mathcal{D}}\!\left[\left(y_0 - \hat{f}(\mathbf{x}_0)\right)^2\right]
=
\underbrace{\sigma^2_{\varepsilon}}_{\text{irreducible}}
+
\underbrace{\left(\mathbb{E}_{\mathcal{D}}[\hat{f}(\mathbf{x}_0)] - f(\mathbf{x}_0)\right)^2}_{\text{bias}^2}
+
\underbrace{\mathbb{E}_{\mathcal{D}}\!\left[\left(\hat{f}(\mathbf{x}_0) - \mathbb{E}_{\mathcal{D}}[\hat{f}(\mathbf{x}_0)]\right)^2\right]}_{\text{variance}} .
\end{equation}
A flexible learner drives the bias term down but leaves the variance term
governed by how much independent information $\mathcal{D}$ carries near
$\mathbf{x}_0$. The key word is \emph{independent}: the variance term does
not shrink when the same underlying relationship is resampled, only when
genuinely new information is added. This is the precise sense in which
synthetic expansion can fail to help, and it is the theoretical statement
of the empirical result in \Cref{ch:discovery}.

\section{The trade-off that synthetic expansion actually makes}

Synthetic generation from literature-constrained rules addresses coverage:
it populates combinations of temperature, concentration, packaging and
indicator that no experimental programme enumerated, in a way that is
consistent with the relationships the source literature reports. That is a
legitimate and common strategy, and it makes the six-specialist
architecture trainable at all.

What it does not do is add independent scientific information. This
distinction deserves a formal statement, because it is the intellectual
core of the entire investigation.

\begin{findingbox}[title=Sample size is not information]
Let $n$ be the number of rows in a training set and let $G$ be the number
of distinct generative mechanisms that produced them. Classical learning
theory bounds depend on the number of \emph{independent} draws. When each
mechanism $g$ contributes $n_g$ rows that are deterministic or
near-deterministic functions of $g$'s parameters, the effective number of
independent units is closer to $G$ than to $n = \sum_g n_g$.

In the V7 corpus, $n = 34\,000$ but $G \in [7, 12]$ per specialist. The
ratio $n/G$ exceeds $1\,000$ for the largest rules. A model that learns
$G$ relationships well will score highly on any test set drawn from those
same $G$ relationships, at any value of $n$, and \emph{regardless of its
own capacity}.
\end{findingbox}

\fig{fig05_rule_counts}{0.92}{%
  The quantity that governs generalization is not the row count but the
  rule count. Each specialist is built from between seven and twelve
  distinct \texttt{source\_rule\_id} values. Counted directly from the
  three V7 specialist files.}

\fig{fig07_rule_concentration}{0.9}{%
  Cumulative share of each specialist's rows as generation rules are added
  largest-first. In four of the six specialists, two rules account for more
  than $90\%$ of all rows.}

\section{Alternatives that existed, and why they are not free}

The panel may reasonably ask whether synthetic expansion was the only
option. It was not, and it is worth being explicit about the alternatives
and their own costs:

\begin{description}[leftmargin=1.5em, style=nextline, font=\bfseries\color{dgreen}]
\item[Train simpler models on the real data only] Fewer parameters would
  reduce variance, but the six-specialist architecture would still divide
  a few hundred observations by six, and the categorical inputs
  (product, indicator, packaging) would still be sparsely populated. This
  option trades one failure mode for another rather than removing it.
\item[Hierarchical or multi-task modelling] Fitting one model with
  category-level random effects, rather than six independent models, would
  pool statistical strength across categories while still allowing
  category-specific behaviour. This is a genuine and largely untried
  alternative in this project, and \Cref{ch:proposals} lists it as a
  recommendation.
\item[Transfer from a mechanistic baseline] Predictive microbiology
  supplies parametric growth models whose parameters could be fitted to the
  small real corpus, with machine learning used only for the residual.
  This constrains the model with physics instead of with synthetic rows.
\item[Collecting more experimental data] The scientifically strongest
  option and the slowest: shelf-life studies run for the duration of the
  shelf life being measured, across multiple temperatures, with
  microbiological enumeration at each time point.
\end{description}

\begin{questionbox}[title=Likely question --- ``Why did you use synthetic data at all?'']
\emph{Suggested answer.} Because the prediction problem is split across
six specialists and a combinatorial input space, and the experimental
corpus available at the time could not populate that space --- some
specialists would have received only tens of rows. Synthetic generation
under literature-constrained rules was used to make the architecture
trainable, not because synthetic rows were believed equivalent to
measurements.

What we have since established, and what this report documents, is that
expansion of this kind increases sample size without increasing
independent information: the corpus has 34\,000 rows but only seven to
twelve generation rules per specialist. That is exactly why the original
evaluation was optimistic, and it is why we now evaluate by holding out
whole generation rules. I would also say plainly that the original
experimental observations are not present in this repository, so I cannot
give you a verified count for them.
\end{questionbox}
